\documentclass[10pt,a4paper]{amsart}
\usepackage[margin=19mm]{geometry}
\usepackage{amsmath,amssymb,mathtools}
\usepackage[scr=rsfs,cal=euler]{mathalfa}
\usepackage{xfrac}
\usepackage[unicode,hidelinks,bookmarks=false]{hyperref}
\newcommand{\ie}{i.e.\ }
\newcommand{\cf}{cf.\ }
\newcommand{\resp}{respectively\ }
\newcommand{\Subsection}{\subsection}
\providecommand{\nobreakdash}{\nobreak\hbox{-}}
\setlength{\emergencystretch}{2em}
\multlinegap0pt
\usepackage{xcolor}
\usepackage{amssymb}
\usepackage{csquotes}
\usepackage[shortlabels]{enumitem}
\newtheorem{lemma}{Lemma}
\newcommand{\R}{\mathbb{R}}
\newcommand{\comm}[2]{%
  {\color{red}#1: #2}}
\title{The unsolvability of the homeomorphism problem}
\author{Stefan Friedl, Tobias Hirsch, Marc Kegel}
\date{Source: arXiv:2603.23630v1 (CC BY 4.0)}
\begin{document}
\maketitle

\noindent\emph{Source and licence.} The unsolvability of the homeomorphism problem. Version arXiv:2603.23630v1 (CC BY 4.0), distributed by arXiv under the Creative Commons Attribution 4.0 International licence, http://creativecommons.org/licenses/by/4.0/, as recorded in the arXiv licence metadata for this exact version and shown on the abstract page https://arxiv.org/abs/2603.23630v1. Full licence notice: \texttt{licenses/arxiv-2603.23630-CC-BY-4.0.txt}.\par\smallskip

\noindent\textit{Editorial scope.} Counted unit: the article's lemma lem:MP-trivtangentbundle (source0.tex lines 916--919). The article's complete original proof of it is delivered with it (source0.tex lines 921--935, one \texttt{proof} environment of 15 lines). No mathematics is written, repaired or re-derived: the statement and the proof are the article's own lines, in the article's own notation, and no retained line is substituted or re-flowed.  No further source-local context is needed: every cross-reference of every retained line resolves inside the two delivered spans.  Nothing was omitted from any retained span, so the package carries no omission record.  No retained line contains a citation, so the wrapper carries no bibliography and no bibliography entry.  No retained line contains a figure environment, an image-inclusion command or drawing code, so no image is delivered and the package declares no asset dependency.  Editorial formatting only, in addition: an AMS \texttt{amsart} preamble that restates the article's own declarations and macros used by the retained lines, namely the environments \texttt{lemma} on separate counters, together with the standard packages those lines need.  The retained environments print in this document as Lemma 1 in order of appearance, whereas the article prints the same items with the numbering of its own full document.  The complete native source of the article as distributed in the arXiv e-print for this exact version is bundled unchanged as \texttt{sources/197092/source0.tex}.
\begin{lemma}
\label{lem:MP-trivtangentbundle}
Among the $2^l$ choices in the above definition, there is at least one choice that leads to a smooth manifold with trivial tangent bundle.
\end{lemma}

\begin{proof}
It suffices to show that in the above construction at least one choice leads to a codimension-zero submanifold of  $\R^5$.
We attach the 1-handles in the standard way to obtain a codimension-zero submanifold $W_1$ of $\R^5$. Let $X_1$ be  the complement of the interior of $W_1$ viewed as a submanifold of $S^5=\R^5\cup \{\infty\}$. Note that $X_1$ admits a handle decomposition with one 0-handle and $k$ 3-handles. In particular, we see that $X_1$, and thus also $X_1\cap \R^5$ is simply connected.

\comm{MK}{Maybe we can make the above a bit more precise. Something like: For every $1$-handle $h^i_1$ of $W_1$ we attach a $2$-handle $h_2^i$ to $W_1$ such that the attaching sphere of the $2$-handle $h_2^i$ intersects the belt sphere of the $1$-handle $h_1^i$ transversely in a single point. Thus the handles cancel and we get a handle decomposition of $D^5$ to which we can glue a $5$-handle to obtain $S^5$. This describes an embedding of $W_1$ into $S^5$. By considering the dual handle decomposition, the exterior of $W_1$ in $S^5$ can be obtained by attaching $k$ $3$-handles to a $0$-handle and thus is simply connected. By seeing $S^5$ as $\R^5\cup \{\infty\}$, we also get an embedding of $W_1$ into $\R^5$ with simply connected exterior $X_1$.}

Now let $C_1,\dots,C_l$ in $\partial W_1$ be smoothly embedded curves that correspond to $r_1,\dots,r_l$. Since 
$X_1\cap \R^5$ is simply connected and since $\dim(X_1)=5>2\cdot 2$
we see that there exist proper smooth embeddings $\varphi_1,\dots,\varphi_l\colon D^2 \to X_1\cap \R^5 $
such that for each $i$ we have $\varphi_i(S^1)=C_i$.
For the proper submanifolds $\varphi_i(D^2)\subseteq X_1$, $i=1,\dots,l$ we pick tubular neighbourhoods. Since $D^2$ and $\R^5$ are orientable, we see that the tubular neighbourhoods are trivial. We now use these trivial tubular neighbourhoods to attach the $l$ 2-handles
and to form $W_2$, which \comm{MK}{corresponds to one of the $2^l$ choices and} is now by construction a codimension-zero submanifold of~$\R^5$.

Finally the $k$ remaining 2-handles are attached in the standard \comm{MK}{I am not sure what standard way means. Maybe we better just say that the $2$-handle attachment from point (4) can be done in $\R^5$.} way and this can evidently be done in $\R^5$.
\end{proof}

\end{document}
